\section{مقدمة}
يُعد التفاعل البشري-الروبوتي (\textenglish{Human-Robot Interaction}) من أهم المجالات التطبيقية في هندسة الأتمتة الحديثة، حيث يسعى المصممون باستمرار إلى إيجاد واجهات تخاطب أكثر مرونة وبديهية من واجهات التحكم التقليدية. في العديد من التطبيقات الصناعية والخدمية، قد يكون المشغل مشغول اليدين أو بعيداً عن لوحة التحكم، مما يجعل توجيه الذراع الروبوتية عبر الأوامر الصوتية حلاً مثالياً لتعزيز الكفاءة والمرونة التشغيلية.

يأتي هذا الفصل استكمالاً للبنية الهرمية للتحكم التي تم تصميمها في الفصل السابق، حيث يتم دمج نظام للتعرف على الكلام (\textenglish{Speech Recognition}) ليعمل كطبقة عليا مسؤولة عن توليد الأهداف الحركية. يهدف هذا النظام إلى التقاط الأوامر الصوتية للمشغل، وتحليلها في الزمن الحقيقي، ومن ثم ترجمتها إلى إحداثيات كارتزية ($X_{\text{ref}}$) أو أوامر تشغيلية للماسك (\textenglish{Gripper})، ليتم تمريرها بعد ذلك إلى طبقة التحكم عالي المستوى (\textenglish{HLC}) لتنفيذ المسار المطلوب.